\section{Conclusion}
\label{conclusion}

This work introduced \sysname, a hypothesis-driven AutoResearch harness for EDA-native GNN exploration. CircuitModelPilot is motivated by a limitation of both fixed-space GNN-NAS and unconstrained LLM-based code exploration. Fixed-space methods can optimize only among modeling choices specified in advance, whereas unconstrained agents can generate broad code changes without ensuring that a candidate corresponds to a meaningful EDA mechanism or that its measured effect is attributable to the claimed intervention.

CircuitModelPilot addresses this gap by making a falsifiable EDA modeling hypothesis the unit of exploration. Semantic co-design provides a grounding vocabulary over representation, computation, and learning signal, allowing the system to formulate mechanisms in terms of EDA objects, dependencies, and analysis rules rather than generic architecture tuning. The framework then realizes each hypothesis as a trace-local candidate and records it in a hypothesis tree. A fixed ResearchContext, scoped editable surfaces, one independent Candidate Review, and fingerprint checks connect the hypothesis to its generated artifacts and comparable execution conditions. The resulting design trace preserves the parent lineage, implementation delta, validation and fixed-test evidence, and analysis outcome, allowing later exploration to extend, ablate, prune, restart, or stop a research direction on the basis of recorded evidence.

An additional principle of CircuitModelPilot is to separate parameter training from trace evaluation. Each run selects a checkpoint on validation and then evaluates that checkpoint on a fixed test split. The resulting validation-and-test evidence informs candidate comparison, analysis, and later research decisions, while neither metric participates in parameter training. This boundary, together with explicit outcome categories of supported, refuted, and inconclusive, treats negative and unstable results as useful research memory rather than discarding them as failed trials.

The evaluation protocol is designed to assess CircuitModelPilot across HLS resource prediction, netlist-level timing prediction, and RTL-to-netlist timing prediction. It examines whether CircuitModelPilot discovers EDA-native mechanisms beyond fixed search spaces, whether the resulting designs improve predictive effectiveness under matched budgets, and whether the trace protocol improves the validity and efficiency of automated exploration. The final quantitative comparison is reported after all baselines, confirmation runs, and matched search trajectories are completed.

Several directions remain open. First, richer domain profiles, diagnostic libraries, and task-specific semantic templates can improve the quality and coverage of candidate hypotheses. Second, stronger confirmation procedures and broader cross-design evaluation can further characterize the stability of discovered mechanisms. Finally, the same trace-grounded methodology could be extended from graph prediction to other EDA learning tasks in which useful modeling choices depend on tool semantics and structured design constraints. More broadly, CircuitModelPilot suggests that reliable automated research requires not only open-ended model generation, but also explicit hypotheses, controlled interventions, and evidence that remains traceable throughout the research loop.
